\chapter{Culture et prise de décision}\label{ch:fm:culture-and-decision-making}

Dans un tournoi de prévision géopolitique de deux ans, où cinq groupes de recherche universitaires rivalisaient pour produire les probabilités les plus justes, un groupe a constaté que trois interventions rendaient ses prévisionnistes meilleurs~:
les former aux probabilités, les laisser partager leurs informations et discuter leurs raisons en équipe, et réunir les meilleurs d'entre eux. Chacune améliorait à la fois la concordance entre les probabilités des prévisionnistes et les issues réelles~;
la netteté avec laquelle ils séparaient les événements qui se produisaient de ceux qui ne se produisaient pas, combinée à une agrégation statistique, la démarche a donné les meilleures prévisions deux années de suite.
Écrire un chiffre, être noté dessus et argumenter au grand jour~:
voilà l'essentiel de ce qu'est une culture du pari,
et une société de trading est un lieu où chaque décision est un pari.

\section{Cultures du pari~: des probabilités, non des opinions}

Un desk de trading tranche en permanence dans l'incertitude~: un signal est-il réel, une stratégie est-elle cassée (chapitre 8), faut-il embaucher, faut-il entrer sur un
marché (chapitre 24), faut-il relever une limite (chapitre 12)~? Une culture qui formule ces croyances en probabilités peut les noter ensuite~; une culture qui les formule en opinions («~je pense que ça marchera~») ne le peut pas,
et n'apprend que des résultats, qui mêlent compétence et chance (chapitre 10).

\begin{definition}[Journal de décisions]\label{def:fm:culture-and-decision-making:journal}
Un \emph{journal de décisions}\index{journal de décisions} consigne chaque décision importante au moment où elle est prise~: les options envisagées, les probabilités
et les espérances de gain attribuées à leurs issues, les raisons, la personne qui a décidé et la date~; le résultat est ajouté plus tard, de sorte que les décisions puissent être
notées sur la base de l'information disponible à l'époque.
\end{definition}

\section{Noter les prévisions et les décisions}

\begin{definition}[Score de Brier]\label{def:fm:culture-and-decision-making:brier}
Le \emph{score de Brier}\index{score de Brier} de prévisions de probabilités $p_i$ d'issues binaires $y_i\in\{0,1\}$ est l'erreur quadratique moyenne
$\frac1N\sum_i(p_i-y_i)^2$~: nul pour des prévisions parfaites, égal à 0.25 pour une prévision constante de un demi, et d'autant meilleur qu'il est bas.
\end{definition}

La décomposition de Murphy (1973) découpe le score en trois parties~: l'incertitude propre aux événements elles-mêmes, $\bar o(1-\bar o)$, que nul prévisionniste ne
contrôle~; la fiabilité, l'écart quadratique moyen entre les prévisions de chaque classe de probabilité et la fréquence à laquelle ces événements se sont produits,
que l'entraînement à la calibration réduit~; et la résolution, la dispersion des fréquences observées par classe autour de la fréquence globale, que la connaissance
augmente. Le score vaut fiabilité moins résolution plus incertitude (\cref{lst:fm:culture-and-decision-making:score})~; la courbe de calibration du livre~7
représente le terme de fiabilité.

\omcode{code/firm/decisionlog/firm_decisionlog.py}{30}{57}{Le score de Brier, les classes de calibration et la décomposition de Murphy.}\label{lst:fm:culture-and-decision-making:score}

\begin{definition}[Agrégation de prévisions]\label{def:fm:culture-and-decision-making:agg}
L'\emph{agrégation de prévisions}\index{agrégation de prévisions} combine les probabilités de plusieurs prévisionnistes pour un même événement en une seule~: en moyennant les
probabilités, en moyennant leurs log-cotes, ou en moyennant les log-cotes puis en les multipliant par un facteur supérieur à un (extrémisation) pour corriger la modération
que produit la moyenne lorsque les prévisionnistes détiennent des informations différentes.
\end{definition}

\begin{proposition}[Cinq avis ou un seul]\label{prop:fm:culture-and-decision-making:rho}
Si les erreurs en log-cote de $K$ prévisionnistes ont chacune pour variance $\sigma^2$ et pour corrélation deux à deux $\rho$, l'erreur de leur moyenne a pour variance
\[
\sigma^2\Bigl(\rho+\frac{1-\rho}{K}\Bigr),
\]
qui décroît de $\sigma^2$ pour un prévisionniste à $\sigma^2/K$ pour des prévisionnistes indépendants, et reste égale à $\sigma^2$ quel que soit $K$ lorsque $\rho=1$.
\end{proposition}

\begin{proof}
$\mathrm{Var}\bigl(\frac1K\sum_ke_k\bigr)=\frac1{K^2}\bigl(K\sigma^2+K(K-1)\rho\sigma^2\bigr)=\sigma^2\bigl(\rho+\frac{1-\rho}K\bigr)$.
\end{proof}

Une discussion qui précède la prévision met en corrélation les erreurs~: les gens convergent vers l'argument le plus bruyant, le premier chiffre énoncé ou le point de vue de la personne
la plus haut placée. La valeur de cinq avis tient à leur indépendance, et c'est pourquoi une bonne pratique de prévision recueille les chiffres avant la discussion
et la discussion avant un second tour.

\section{Tutoriel~: cinq avis ou un seul}\label{tut:fm:culture-and-decision-making:tut}

\begin{tutorial}
\textbf{Objectif.} Noter une année de prévisions d'un desk, décomposer les scores, et comparer l'agrégation de cinq prévisions indépendantes avec l'enregistrement d'un consensus unique
obtenu après discussion. \textbf{État final~:} les courbes de calibration (\cref{fig:fm:culture-and-decision-making:calib}) et les \omterm{def:fm:culture-and-decision-making:brier}{scores de Brier} par règle
(\cref{fig:fm:culture-and-decision-making:rho}).
\begin{tutsteps}
\item \textbf{Les événements.} 400 événements par an, avec des probabilités vraies tirées d'une loi Beta(2, 2)~; la vérité elle-même obtient un score de 0.214, et une constante
de un demi 0.25 (\texttt{fm\_culture.year}, synthétique).
\item \textbf{Les prévisionnistes.} Cinq voient la log-cote vraie avec un bruit d'écart type 0.8, et étirent leur log-cote d'un facteur de 1.0 à 1.6~: tous
sauf le premier sont trop confiants.
\item \textbf{Les scores.} \texttt{firm.decisionlog.brier}, \texttt{murphy} et \texttt{calibration}~; \texttt{aggregate} de trois façons.
\item \textbf{La discussion.} \texttt{fm\_culture.five\_or\_one} met en corrélation les erreurs des prévisionnistes à $\rho$ allant de 0 à 1 et note le consensus
sur vingt années simulées.
\end{tutsteps}
\end{tutorial}

\begin{omfigure}
\begin{tikzpicture}
\begin{axis}[width=8cm,height=7cm,xlabel={\small probabilité prévue},ylabel={\small fréquence observée},xmin=0,xmax=1,ymin=0,ymax=1,
  tick label style={font=\footnotesize},grid=major,grid style={gray!20},table/col sep=comma,
  legend cell align=left,legend style={font=\scriptsize,at={(1.05,0.5)},anchor=west}]
\addplot[gray,dashed] coordinates {(0,0) (1,1)};
\addplot[omThm,thick,mark=square*,mark size=1.3pt] table[x=forecast,y=observed]{figdata/desk/26-culture-and-decision-making/calib_f1.csv};
\addplot[omProp,thick,mark=triangle*,mark size=1.6pt] table[x=forecast,y=observed]{figdata/desk/26-culture-and-decision-making/calib_f5.csv};
\addplot[omDef,very thick,mark=*,mark size=1.3pt] table[x=forecast,y=observed]{figdata/desk/26-culture-and-decision-making/calib_agg.csv};
\legend{calibration parfaite,prévisionniste 1,prévisionniste 5 (trop confiant),moyenne des log-cotes}
\end{axis}
\end{tikzpicture}

\omcaption{Courbes de calibration pour une année de 400 prévisions~: le prévisionniste 5, qui étire ses log-cotes d'un facteur 1.6, attribue des probabilités élevées à des événements qui se produisent
bien moins souvent~; la moyenne des cinq prévisionnistes est proche de la diagonale. Données~: \texttt{firm.decisionlog.calibration}.}\label{fig:fm:culture-and-decision-making:calib}
\end{omfigure}

Sur l'année du chapitre, les cinq prévisionnistes obtiennent des scores compris entre 0.225 et 0.262, contre 0.214 pour la vérité elle-même. Le plus confiant, le prévisionniste 5,
présente un terme de fiabilité de 0.040 contre 0.017 pour le prévisionniste 1, avec une résolution à peu près identique (0.026 et 0.027)~: son erreur supplémentaire tient entièrement
à une mauvaise calibration. La moyenne des cinq probabilités obtient un score de 0.220 et la moyenne de leurs log-cotes 0.223, toutes deux meilleures que chacun des prévisionnistes
(\cref{fig:fm:culture-and-decision-making:brier})~; le terme de fiabilité de la moyenne
tombe à 0.010 et sa résolution monte à 0.039, puisque les erreurs indépendantes se compensent en partie. L'extrémisation de la moyenne l'aggrave
(0.237)~: ces prévisionnistes sont déjà trop confiants, et l'extrémisation n'aide que les prévisionnistes dont la moyenne est trop timide.

\begin{omfigure}
\begin{tikzpicture}
\begin{axis}[width=10.5cm,height=5.4cm,xbar,bar width=7pt,xmin=0.20,xmax=0.27,xlabel={\small score de Brier (plus bas est meilleur)},ytick={0,1,2,3,4,5,6,7,8},
  yticklabels={prévisionniste 1,prévisionniste 2,prévisionniste 3,prévisionniste 4,prévisionniste 5,moyenne,moyenne des log-cotes,extrémisée,vérité},y dir=reverse,ymin=-0.6,ymax=8.6,
  tick label style={font=\footnotesize},x tick label style={/pgf/number format/.cd,fixed,precision=2},grid=major,grid style={gray!20},
  table/col sep=comma,nodes near coords,nodes near coords style={font=\scriptsize,/pgf/number format/.cd,fixed,precision=3},point meta=x]
\addplot[fill=omDef!55,draw=omDef] table[y=pos,x=brier]{figdata/desk/26-culture-and-decision-making/brier.csv};
\end{axis}
\end{tikzpicture}

\omcaption{\omterm{def:fm:culture-and-decision-making:brier}{Scores de Brier} des 400 prévisions d'une année~: cinq prévisionnistes, trois façons de les agréger, et les probabilités propres de la vérité, le plancher
qu'aucun prévisionniste ne peut battre en moyenne. Données~: \texttt{fm\_culture.scores}.}\label{fig:fm:culture-and-decision-making:brier}
\end{omfigure}

\begin{omfigure}
\begin{tikzpicture}
\begin{axis}[width=10.5cm,height=5cm,xlabel={\small corrélation des erreurs après discussion, $\rho$},ylabel={\small score de Brier},xmin=0,xmax=1,ymin=0.20,ymax=0.245,
  tick label style={font=\footnotesize,/pgf/number format/.cd,fixed,precision=3},grid=major,grid style={gray!20},table/col sep=comma,
  legend cell align=left,legend style={font=\scriptsize,at={(0.5,-0.3)},anchor=north,/tikz/every even column/.append style={column sep=8pt}},legend columns=2]
\addplot[omDef,very thick,mark=*,mark size=1.3pt] table[x=rho,y=one]{figdata/desk/26-culture-and-decision-making/rho.csv};
\addplot[omProp,thick,dashed] table[x=rho,y=independent]{figdata/desk/26-culture-and-decision-making/rho.csv};
\legend{un seul consensus après discussion,cinq prévisions indépendantes}
\end{axis}
\end{tikzpicture}

\omcaption{\omterm{def:fm:culture-and-decision-making:brier}{Score de Brier} de la prévision du desk, moyenné sur vingt années simulées~: cinq prévisions indépendantes agrégées, face à un
consensus dont une discussion a corrélé les erreurs à $\rho$. Données~: \texttt{fm\_culture.five\_or\_one}.}\label{fig:fm:culture-and-decision-making:rho}
\end{omfigure}

Sur vingt années simulées, cinq prévisions indépendantes obtiennent en moyenne un \omterm{def:fm:culture-and-decision-making:brier}{score de Brier} de 0.207~; un consensus issu d'une discussion qui corrèle les erreurs à
0.5 obtient un score de 0.224, et à 1 (chacun répète la même opinion), 0.237 (\cref{prop:fm:culture-and-decision-making:rho}). La différence, 0.017 pour
$\rho=0.5$, représente plus de la moitié de l'écart entre le score moyen des prévisionnistes de l'année du chapitre et celui de la vérité~: la discussion détruit la plus grande part de ce que valaient les
prévisions indépendantes.

\section{Post-mortems et pré-mortems}

\begin{definition}[Biais du résultat]\label{def:fm:culture-and-decision-making:outcome}
Le \emph{biais du résultat}\index{biais du résultat} (outcome bias) consiste à juger la qualité d'une décision par son résultat plutôt que par l'information et le raisonnement disponibles
au moment où elle a été prise, de sorte que les bonnes décisions qui ont perdu sont qualifiées d'erreurs et que les mauvaises décisions qui ont gagné sont récompensées.
\end{definition}

Baron et Hershey (1988) ont étudié ce biais expérimentalement~; le \omterm{def:fm:culture-and-decision-making:journal}{journal de décisions} en est l'antidote. Dans l'année du chapitre, le prévisionniste 2, le meilleur, parie à
cote égale sur les 362 événements qu'il évalue au-dessus de 55~\% ou en dessous de 45~\%~: chaque pari a une espérance positive selon sa propre prévision, et 123 d'entre eux, 34~\%,
perdent. Un relecteur qui juge au résultat qualifie de mauvais le tiers des bonnes décisions du desk. Le même journal révèle le vrai défaut du prévisionniste 2~: il gagne
66~\% de ses paris alors que ses prévisions s'établissent en moyenne à 77.7~\%, un excès de confiance que seul un registre des probabilités énoncées peut révéler.

\begin{definition}[Pré-mortem, équipe rouge]\label{def:fm:culture-and-decision-making:premortem}
Un \emph{pré-mortem}\index{pré-mortem} est une revue tenue avant qu'une décision soit mise en œuvre, au cours de laquelle l'équipe imagine que le projet a échoué et
note pourquoi, pour faire remonter les risques que l'optimisme dissimule. Une \emph{équipe rouge}\index{equipe rouge@équipe rouge} (red team) est un groupe chargé de contester un plan ou d'attaquer un
système, indépendamment de ceux qui l'ont conçu.
\end{definition}

Les post-mortems (les revues sans recherche de coupable du livre~15, blameless) portent sur un résultat passé~; les \omterm{def:fm:culture-and-decision-making:premortem}{pré-mortems} portent sur une décision à venir. Les deux ne fonctionnent que s'ils séparent la
qualité de la décision de la part de hasard du résultat, ce qui suppose le journal.

\section{Comment les bonnes sociétés gèrent leurs désaccords}

Une société gère bien ses désaccords lorsqu'elle recueille les avis avant de les discuter, les exprime sous forme de nombres, laisse la prévision d'un junior compter autant que celle
d'un senior, désigne quelqu'un pour défendre le point de vue adverse, et évalue chacun a posteriori. La revue de recherche et le pré-enregistrement du livre~7 sont la même
discipline appliquée à la recherche~; le journal de recherche est un \omterm{def:fm:culture-and-decision-making:journal}{journal de décisions} pour les stratégies.

\begin{method}[Faire vivre une culture du pari]\label{met:fm:culture-and-decision-making:run}
\begin{enumerate}
\item Tenir un \omterm{def:fm:culture-and-decision-making:journal}{journal de décisions} pour chaque décision importante, avec les probabilités et les espérances écrites avant le résultat.
\item Recueillir les prévisions de façon indépendante avant les réunions~; discuter~; recueillir à nouveau~; agréger les log-cotes.
\item Noter les prévisionnistes chaque trimestre avec le \omterm{def:fm:culture-and-decision-making:brier}{score de Brier} et sa décomposition~; travailler la calibration là où le terme de fiabilité est élevé.
\item Revoir les décisions d'après le processus et l'information disponible, et non d'après le résultat~; tenir des \omterm{def:fm:culture-and-decision-making:premortem}{pré-mortems} pour les grandes décisions et donner à une \omterm{def:fm:culture-and-decision-making:premortem}{équipe rouge}
le droit d'être entendue.
\item Récompenser la bonne calibration et le bon processus, et non les résultats chanceux (chapitre~10).
\end{enumerate}
\end{method}

\section{Construction~: le journal de décisions}\label{bld:fm:culture-and-decision-making:decisionlog}

\begin{build}
\bfield{Rôle} Un \omterm{def:fm:culture-and-decision-making:journal}{journal de décisions} et la notation de ses prévisions~: \omterm{def:fm:culture-and-decision-making:brier}{scores de Brier}, décomposition de Murphy, calibration, agrégation et une
vérification du \omterm{def:fm:culture-and-decision-making:outcome}{biais du résultat}.
\bfield{Interface} \texttt{firm.decisionlog}~: \texttt{Entry}, \texttt{brier}, \texttt{murphy}, \texttt{calibration}, \texttt{aggregate},
\texttt{logit}, \texttt{expit}, \texttt{outcome\_bias}.
\bfield{Règles} Les probabilités sont enregistrées avant les résultats~; la décomposition est exacte lorsque les prévisions sont constantes à l'intérieur des classes~; l'agrégation porte sur
les mêmes événements.
\bfield{Tests d'acceptation} \texttt{code/firm/decisionlog/tests/}~: l'identité de la décomposition sur les prévisions regroupées par classes~; les trois agrégations sur des
nombres calculés à la main~; le décompte du \omterm{def:fm:culture-and-decision-making:outcome}{biais du résultat}.
\bfield{Pour aller plus loin} Noter des prévisions continues~; pondérer les prévisionnistes selon leur compétence passée~; un \omterm{def:fm:culture-and-decision-making:journal}{journal des décisions} de limites et d'embauche lié au
registre du chapitre~12.
\end{build}

\begin{omsources}
\item B.~A.~Mellers et al., «~Psychological strategies for winning a geopolitical forecasting tournament~», \emph{Psychological Science} 25(5), 2014.
\item G.~W.~Brier, \emph{Monthly Weather Review} 78(1), 1950; A.~H.~Murphy, \emph{Journal of Applied Meteorology} 12, 1973; V.~A.~Satopaa et al.,
\emph{International Journal of Forecasting}, 2014.
\item J.~Baron et J.~C.~Hershey, «~Outcome bias in decision evaluation~», \emph{Journal of Personality and Social Psychology} 54(4), 1988.
\end{omsources}

\section{Exercices}

\begin{exercise}[$\star$]\label{exo:fm:culture-and-decision-making:1}
Calculer le \omterm{def:fm:culture-and-decision-making:brier}{score de Brier} de la prévision 0.7 pour trois événements dont deux se sont réalisés. Quel score obtient une prévision constante de 0.5~?
\end{exercise}

\begin{exercise}[$\star$]\label{exo:fm:culture-and-decision-making:2}
Calculer la moyenne des probabilités 0.6 et 0.8, ainsi que celle de leurs log-cotes. Laquelle est la plus éloignée de un demi~?
\end{exercise}

\begin{exercise}[$\star$]\label{exo:fm:culture-and-decision-making:3}
Avec cinq prévisionnistes dont les erreurs sont indépendantes, calculer de quel facteur la moyenne divise la variance de l'erreur, puis refaire le calcul pour une corrélation de 0.5.
\end{exercise}

\begin{exercise}[$\star\star$]\label{exo:fm:culture-and-decision-making:4}
Expliquer pourquoi le prévisionniste 5 obtient un score moins bon que le prévisionniste 1 alors que sa résolution est à peu près la même.
\end{exercise}

\begin{exercise}[$\star\star$]\label{exo:fm:culture-and-decision-making:5}
Expliquer pourquoi l'extrémisation nuit dans l'année du chapitre, et dans quels cas elle aiderait.
\end{exercise}

\begin{exercise}[$\star\star$]\label{exo:fm:culture-and-decision-making:6}
Expliquer comment un desk doit conduire une réunion pour décider si une stratégie est défaillante, compte tenu de la \cref{prop:fm:culture-and-decision-making:rho}.
\end{exercise}

\begin{exercise}[$\star\star\star$]\label{exo:fm:culture-and-decision-making:7}
\emph{Programmation.} Noter les paris du journal selon le résultat et selon le processus. Quelle part des décisions d'espérance positive une revue fondée sur le résultat qualifierait
d'erreurs~?
\end{exercise}

\begin{exercise}[$\star\star\star$]\label{exo:fm:culture-and-decision-making:8}
\emph{Trouver la faille.} «~On a gagné de l'argent sur cette transaction, donc c'était une bonne décision.~»
\end{exercise}

\section{Problème~: Cinq avis ou un seul}

\begin{problem}[{Problème du week-end --- cinq avis ou un seul}]\label{pb:fm:culture-and-decision-making:1}
Un responsable de la recherche veut remplacer la réunion hebdomadaire de consensus du desk par un processus de prévision noté.

\medskip\noindent\textbf{Partie I --- Principes.}
\begin{enumerate}
\item Que révèle le tournoi de prévision de l'accroche~?
\item Définir un \omterm{def:fm:culture-and-decision-making:journal}{journal de décisions}.
\item Définir le \omterm{def:fm:culture-and-decision-making:brier}{score de Brier} et donner sa valeur pour une prévision constante égale à un demi.
\item Énoncer la décomposition de Murphy et ce que mesure chacun de ses termes.
\end{enumerate}

\medskip\noindent\textbf{Partie II --- L'année.}
\begin{enumerate}[resume]
\item Donner les scores des cinq prévisionnistes et celui de l'issue réelle.
\item Décomposer les scores des prévisionnistes 1 et 5 et celui de la moyenne.
\item Comparer les trois règles d'agrégation.
\item Lire les courbes de calibration.
\end{enumerate}

\medskip\noindent\textbf{Partie III --- La discussion.}
\begin{enumerate}[resume]
\item Définir l'\omterm{def:fm:culture-and-decision-making:agg}{agrégation de prévisions}.
\item Énoncer et démontrer la \cref{prop:fm:culture-and-decision-making:rho}.
\item Donner les \omterm{def:fm:culture-and-decision-making:brier}{scores de Brier} des prévisions indépendantes et d'un consensus pour $\rho=0.5$ et $\rho=1$.
\item Pourquoi la discussion corrèle-t-elle les erreurs~?
\item Expliquer comment conduire la réunion à la place.
\end{enumerate}

\medskip\noindent\textbf{Partie IV --- Culture.}
\begin{enumerate}[resume]
\item Définir le \omterm{def:fm:culture-and-decision-making:outcome}{biais du résultat} et donner la preuve que le journal en fournit.
\item Définir un \omterm{def:fm:culture-and-decision-making:premortem}{pré-mortem} et une \omterm{def:fm:culture-and-decision-making:premortem}{équipe rouge} (red team).
\item Comment la rémunération doit-elle refléter la compétence en prévision~?
\item Quelles informations consigner dans un journal pour une hausse de limite~?
\item Que peut-il mal tourner lorsqu'on note les personnes~?
\item Énoncer le \emph{résultat nommé}~: le gain de \omterm{def:fm:culture-and-decision-making:brier}{score de Brier} obtenu en agrégeant cinq prévisions indépendantes plutôt qu'une seule prévision discutée dont
les erreurs sont corrélées à un niveau donné.
\item En deux phrases, rédiger la proposition.
\end{enumerate}
\end{problem}

\section{Questions d'entretien}

\begin{interviewq}[$\star$ \iqroles{chercheur}]\label{iq:fm:culture-and-decision-making:1}
Qu'est-ce qu'un \omterm{def:fm:culture-and-decision-making:brier}{score de Brier}, et que signifie un score de 0.25~?
\end{interviewq}

\begin{interviewq}[$\star$ \iqroles{trader}]\label{iq:fm:culture-and-decision-making:2}
Vous avez perdu de l'argent sur une transaction que vous jugez toujours juste. Comment faites-vous la différence~?
\end{interviewq}

\begin{interviewq}[$\star\star$ \iqroles{chercheur}]\label{iq:fm:culture-and-decision-making:3}
Pourquoi moyenner les log-cotes plutôt que les probabilités~?
\end{interviewq}

\begin{interviewq}[$\star\star$ \iqroles{chercheur, trader}]\label{iq:fm:culture-and-decision-making:4}
Les erreurs de cinq personnes sont corrélées à 0.5. De combien la moyenne de leurs prévisions réduit-elle la variance de l'erreur~?
\end{interviewq}

\begin{interviewq}[$\star\star$ \iqroles{risque}]\label{iq:fm:culture-and-decision-making:5}
Comment organiseriez-vous un \omterm{def:fm:culture-and-decision-making:premortem}{pré-mortem} pour une nouvelle entrée sur un marché~?
\end{interviewq}

\begin{interviewq}[$\star\star\star$ \iqroles{chercheur}]\label{iq:fm:culture-and-decision-making:6}
Dérivez la décomposition de Murphy du \omterm{def:fm:culture-and-decision-making:brier}{score de Brier}.
\end{interviewq}
